% Generated draft: requires mathematical and source review.
\documentclass[11pt]{article}
\usepackage[margin=25mm]{geometry}
\usepackage{fontspec}
\setmainfont{Latin Modern Roman}
\usepackage{amsmath,amssymb,mathtools,unicode-math}
\setmathfont{Latin Modern Math}
\usepackage{enumitem,needspace,xurl,hyperref}
\setlength{\parindent}{0pt}
\setlength{\parskip}{5pt}
\setlength{\emergencystretch}{3em}
\newcommand{\N}{\mathbb{N}}
\newcommand{\Z}{\mathbb{Z}}
\newcommand{\Q}{\mathbb{Q}}
\newcommand{\R}{\mathbb{R}}
\newcommand{\C}{\mathbb{C}}
\newcommand{\F}{\mathbb{F}}
\newcommand{\A}{\mathbb{A}}
\newcommand{\PP}{\mathbb{P}}
\newcommand{\E}{\mathbb{E}}
\newcommand{\eps}{\varepsilon}
\newcommand{\dd}{\,\mathrm{d}}
\begin{document}
\section*{UnsolvedMath ID: 10600008}
\textbf{Source problem code:} catalogue entry
\subsection{Definitions and mathematical statement}
A diagrammatic biquandle presentation assigns generators to semiarcs and crossing relations to their labels. A switch \(B:X^2\to X^2\) describes a classical crossing (with \(B^{-1}\) for the inverse crossing), and a virtual switch \(V:X^2\to X^2\) describes a virtual crossing. The operations must satisfy the identities implementing all allowed classical, virtual, and mixed Reidemeister moves, including the appropriate nondegeneracy and type-I conditions. The excerpt does not fix a particular convention for these axioms.

\emph{Construction problem.} Construct presentations of a virtual biquandle with a linear, noncommutative representation at classical crossings and interesting structure at virtual crossings. In a switch formulation the classical switch should admit the requested linear realization, for instance as an invertible block matrix on \(X\oplus X\) for a module \(X\) over a noncommutative coefficient ring. A free right-module convention makes left matrix multiplication a right-linear map; the source does not choose a ring or a module side. The virtual operation must be compatible with the classical one and give the intended additional structure.

For example, writing \(B_1=B\times\operatorname{id}\), \(B_2=\operatorname{id}\times B\), and similarly for \(V\), compatibility in a standard braid-switch convention includes
\[
 B_1B_2B_1=B_2B_1B_2,\qquad
 V^2=\operatorname{id},\qquad
 V_1V_2V_1=V_2V_1V_2,\qquad
 V_1V_2B_1=B_2V_1V_2.
\]
These illustrate necessary braid-level relations in that convention, not a replacement for the full biquandle axioms or an asserted solution. A proposed presentation must also verify the remaining move relations and explain what mathematical information its virtual-crossing structure provides. The adjective ``interesting'' supplies no precise nontriviality or distinguishing criterion, and the meaning of ``linear (non-commutative)'' is not specified beyond the source wording.

The source credits Bartholomew and Fenn with an initial computer search producing various biquandle decorations at both classical and virtual crossings. The cited paper's bibliographic details and examples are not supplied.
\subsection{Short English statement}
Build virtual-biquandle presentations with a noncommutative linear rule at classical crossings and meaningful compatible rules at virtual crossings. Full move invariance and the intended significance of the virtual structure must be specified; the source reports preliminary computer-search examples by Bartholomew and Fenn.
\subsection{Sources}
\begin{itemize}
\item Local UnsolvedMath registry, permanent ID 10600008, source code catalogue entry, ``Virtual-knot problem 8 --- Virtual Biquandle.'' No individual proposer is supplied.
\item \emph{Unsolved Problems in Virtual Knot Theory and Combinatorial Knot Theory}, revised survey, problem-list item 8. Supplied URL: \url{https://arxiv.org/abs/1409.2823}; recorded access 2026-07-29; extraction source-tex. The dedicated section contains 60 top-level enumerated items; compound items are retained intact.
\item The statement credits Bartholomew and Fenn's computer search, but the paper citation is missing. The supplied research summary also mentions Bartholomew--Fenn searches and Nelson's biquandle census without full references.
\item Supplied status open and OPEN-TRIAGE. The summary says no published resolution of this specific construction was found; this is an editorial search assessment, not an independently verified fact. Current primary-literature status and redistribution terms for the publicly accessible author/publisher copy remain NEEDS\_REVIEW.
\end{itemize}
\end{document}
